\section{Introduction}
Lexical processing links knowledge about a referent or intended message with a word identity and its phonological form. A computational account must specify both the representations available at these levels and the operations that transform one representation into another. These requirements become particularly important when a model is implemented using spiking neurons: a diagram of connected lexical levels does not by itself explain what the neural populations represent, how a mapping is realized, or what a plotted activation value measures.

Previous work has implemented lexical processing within the Neural Engineering Framework (NEF) and Semantic Pointer Architecture (SPA). \citet{stille2020} investigated the mental lexicon in long-term and working memory and simulated naming with semantic and phonological cues. Related work addressed production--perception interactions and internal monitoring \citep{kroeger2016}, while \citet{kroeger2023} placed the approach within a broader framework for speech production. These studies provide the scientific context for the present work. The present study examines the same basic lexical organization of concepts, lemmas, and phonological forms within more narrowly defined retrieval tasks. The restriction concerns the task setting and the surrounding task architecture, rather than a reduction of the lexical inventory or its three representational levels. The present implementation is not a reproduction of the complete task architecture or the published time courses of those models.

The contribution here is a formal and implementation-level account of the lexical core. Three complementary aspects organize this account. First, our approach links lexical graph representations to their vector encodings and to the representation of those vectors in populations of spiking neurons. This makes it possible to trace how linguistic relations contribute to representational similarity and how the resulting representations are realized in neural activity. Second, we distinguish similarity within a representational level from association between levels: the former describes overlap among representations, while the latter links corresponding concepts, lemmas, and phonological forms for lexical retrieval. Third, we examine activation patterns at the lexical levels and assess which lexical item is retrieved at the output of each processing stage. This relates the similarity structure of the representations to the outcome of lexical processing.

To develop these three aspects, the remainder of this paper proceeds from the theoretical framework to the construction of a 613-entry lexicon and its implementation in separate perception and production pathways. We then present the simulation methods and baseline results for eight selected words from the lexicon, followed by a discussion of their interpretation, the scope of the model, and directions for further investigation. The processing tasks begin with complete phonological form or conceptual vectors. Acoustic recognition, visual object recognition, and articulation are outside the present model. This restricted scope allows the relationships among representational geometry, associative retrieval, and neural realization to be examined directly. The transition from phonological form or syllabic representation to articulatory control is reserved for a separate account.

The primary aim of this paper is to provide a transparent description of the model, supported by illustrative simulations. The resulting simulation corpus provides a basis for broader evaluations, including additional inputs, model realizations, and comparisons with human behavior, as outlined in Section~\ref{sec:limitations}.
